\begin{proof}
(i) A segment of $\bar X_{\mathcal K}$ parallel to $e_i$ maps, for each $t$
with $i\in E_t$, to a segment of $\Pi_t$ parallel to $e_{j(t,i)}$, and
leaves the other value factors constant. Along it, $F_0$ has second
derivative $(H_0)_{ii}$, and $\psi_t+\sigma_{t,j(t,i)}v_i^2/2$ is concave by
Proposition~\ref{prop:vf-curv}. A sum of concave functions is concave. The
checking cost was discussed after Definition~\ref{def:leaf}.

(ii) If $L^+=0$, $V$ is concave along every coordinate segment of
$\bar X_{\mathcal K}$. Replacing coordinates one at a time by a better
endpoint does not increase $V$, so a vertex of $\bar X_{\mathcal K}$, which
lies in $X_{\mathcal K}$, is optimal. Tree dynamic programming over the two
endpoint nodes, with value-factor values from
Lemma~\ref{lem:cv-bits}, finds it; every bag table has at most $2^p$
entries. Private coordinates are recovered by one convex quadratic program
per block.

(iii) \emph{Approximation.} By (i), $V$ has upper coordinate curvatures
$L_i^+$; coordinates with $L_i\le0$ get the two endpoint nodes. Each
$\psi_t$ is assigned to a bag containing $E_t$. An exact convex quadratic
program \cite{KozlovTarasovKhachiyan1980}, or the leaf of a supplied
certificate that contains the query point, gives $\psi_t$ and a minimizer
$y^{(t)}$ exactly in time polynomial in $I$ and the encoding length of the
point (Lemma~\ref{lem:cv-bits}), with a fixed exponent. So
Theorem~\ref{thm:valuefn}(a),(b) apply with $L^V=L^+$; the encoding length
of the $L_i^+$, and hence $\log(1/L^+)$, is polynomial in $I$. There are at
most $rK^p$ value-factor evaluations per stage.

\emph{Exact output.} Lemma~\ref{lem:cv-height}, applied to the integer
slices of the retained box, bounds the denominator of $\OPT$ without
uniqueness, and under growth the denominators of the continuous coordinates
of $v^*$. The exact minimizer is obtained by rational reconstruction of the
incumbent and accepted by Proposition~\ref{prop:accept} with this
denominator bound; Appendix~\ref{app:recourse-convex} gives the procedure
and its analysis.

(iv) $\sigma_t\ge0$ by Proposition~\ref{prop:vf-curv}(b), so
$L_i\le(H_0)_{ii}$. Along $e_i$ the quadratic $F$ has second derivative
$(H_0)_{ii}$, so $L'\ge\max\{(H_0)_{ii},0\}\ge L_i^+$ for every
$i\in\mathcal K$, and $L'\ge L^+$. By Lemma~\ref{lem:valuefunction}(b),
$g$ is also a growth constant of $V$, as (iii) requires, and
$\kappa(L^+,g)\le\kappa(L',g)$.
\end{proof}
